\section{Surrogate Analysis}

\subsection{Recovery Preservation}

\textbf{Lemma 1.}
For a fixed sparse flexible code $B$, any row assignment or row permutation
preserves the algebraic recovery capability of the code.

\textit{Proof sketch.}
Assignment changes when rows arrive, not the row space generated by all rows.
For any permutation matrix $P$, $\mathrm{rank}(B)=\mathrm{rank}(PB)$ and
$\mathbf{1}_m \in \mathrm{span}(B^\top)$ if and only if
$\mathbf{1}_m \in \mathrm{span}((PB)^\top)$.  Therefore scheduling can improve
the first decodable time without weakening the underlying recovery threshold.

\subsection{Decode-Mass Surrogate}

Fix a deadline $u$.  Suppose worker $j$ completes its assigned row pair before
$u$ with probability $p_j(u)$, and $p_j(u)$ is monotone in worker speed.  A
natural surrogate for early decodability is expected completed decode mass:
\begin{equation}
  \label{eq:decode-mass-surrogate}
  \max_\pi \sum_i p_{\pi(i)}(u) R_i .
\end{equation}

\textbf{Lemma 2.}
The optimal assignment for the surrogate in
Equation~\eqref{eq:decode-mass-surrogate} pairs the largest $R_i$ with the
largest $p_j(u)$, the second largest with the second largest, and so on.

\textit{Proof sketch.}
The result follows from the rearrangement inequality.  For any inversion where
$R_a \ge R_b$ but $p_{\pi(a)} < p_{\pi(b)}$, swapping the two assignments
changes the objective by
$(R_a-R_b)(p_{\pi(b)}-p_{\pi(a)}) \ge 0$.  Repeatedly removing inversions gives
the monotone assignment.

\subsection{Mismatch Creates Scheduling Opportunity}

The preceding surrogate also gives a way to state when decode-aware scheduling
has room to help.  Let $q_j(u)=p_j(u)$ be the probability that worker $j$
finishes before deadline $u$, and let $\pi_0$ be the static row-pair placement.
Define the normalized decode-speed mismatch
\begin{equation}
  \label{eq:mismatch-score}
  M(\pi_0,u)
  =
  \frac{
    \max_\pi \sum_i R_i q_{\pi(i)}(u)
    -
    \sum_i R_i q_{\pi_0(i)}(u)
  }{
    \max_\pi \sum_i R_i q_{\pi(i)}(u)
    -
    \min_\pi \sum_i R_i q_{\pi(i)}(u)
  } . 
\end{equation}
In Equation~\eqref{eq:mismatch-score}, the numerator is the decode-mass
headroom left by the static placement
$\pi_0$: it compares the best assignment at deadline $u$ with the assignment
the fixed runtime would have used.  The denominator is the full feasible range
between the best and worst assignments, so the ratio normalizes the score to
$[0,1]$ when worker probabilities and row priorities are not both constant.
$M=0$ means the static placement is already optimal for this surrogate; $M=1$
means it is maximally anti-aligned.  The score is an opportunity signal, not a
latency guarantee.

\textbf{Lemma 3.}
If $M(\pi_0,u)=0$, no reassignment can improve the expected decode-mass
surrogate at deadline $u$.  If $M(\pi_0,u)>0$, the best monotone reassignment
improves the surrogate by exactly $M(\pi_0,u)$ times the surrogate's assignment
range.

\textit{Proof sketch.}
The definition subtracts the static surrogate value from the optimal surrogate
value and normalizes by the feasible range.  Lemma 2 identifies the optimizer
as the monotone assignment.  Therefore zero mismatch means there is no
decode-mass opportunity, while positive mismatch quantifies the available
surrogate improvement.

\subsection{Prefix Reduction Realizes Latency Gain}

The mismatch score says that reassignment has room to help, but it does not by
itself imply lower latency.  The benefit appears only when a decodable row set
arrives earlier.  For an assignment $\pi$, order returned rows by arrival time
and let $T_{\pi,(k)}$ be the time of the $k$th arrival.  Define
\begin{equation}
  \label{eq:first-decodable-prefix}
  K_\pi = \min\{k : \text{the first } k \text{ arrivals under } \pi
  \text{ are decodable}\}.
\end{equation}
Equation~\eqref{eq:first-decodable-prefix} defines the first-decodable prefix
length: the number of completed rows needed before the row-span decoder can
recover the full gradient.  If $h_\pi$ is scheduling and monitoring overhead,
the first-decode time is
\begin{equation}
  \label{eq:first-decode-overhead}
  \tau_\pi = T_{\pi,(K_\pi)} + h_\pi .
\end{equation}

This gives the guard's decision criterion.  Compare static placement $\pi_0$
with adaptive placement $\pi$.  If the candidate does not increase the required
prefix ($K_\pi\le K_{\pi_0}$), the predicted gain can be written as
\begin{equation}
  \label{eq:prefix-gain-decomposition}
  \tau_{\pi_0}-\tau_\pi
  =
  \underbrace{T_{\pi_0,(K_{\pi_0})}-T_{\pi_0,(K_\pi)}}_{\text{fewer rows waited for}}
  +
  \underbrace{T_{\pi_0,(K_\pi)}-T_{\pi,(K_\pi)}}_{\text{same prefix arrives earlier}}
  -
  h_\pi .
\end{equation}
Equation~\eqref{eq:prefix-gain-decomposition} separates three effects: prefix
reduction, faster placement for the same prefix length, and overhead.  This is
why the guard rejects candidates that predict a longer prefix or too much
post-decode work: speed alone is insufficient unless the candidate exposes a
decodable prefix early enough to pay for scheduling, communication, and
cancellation overhead.

\subsection{Surrogate Rationale and Limits}

The exact event of interest is whether the rows completed by deadline $u$ form a
decodable set.  If $\mathcal{D}$ is the family of row subsets that can decode
the target vector, then
\begin{equation}
  \label{eq:decodable-event}
  \Pr[\tau(\pi,B) \le u]
  =
  \Pr[S_\pi(u) \in \mathcal{D}].
\end{equation}
Computing Equation~\eqref{eq:decodable-event} exactly requires reasoning over
many completed
subsets.  The decode-mass surrogate uses a scalar proxy instead: let
$Z_{\pi,i}(u)$ indicate whether row pair $i$ finishes before $u$, and define
$X_\pi(u)=\sum_i R_i Z_{\pi,i}(u)$.  A threshold relaxation treats a row set as
promising when $X_\pi(u)\ge\theta$.  Increasing
$\mathbb{E}[X_\pi(u)]$ therefore improves this relaxed event when variance does
not grow enough to offset the mean gain.

This is a runtime proxy, not a correctness argument.  A large $R_i$ suggests
that row pair $i$ is useful to expose early, but it does not certify that any
returned subset is decodable.  The runtime uses the surrogate only to choose an
assignment; correctness still comes from the exact online row-span decoder.  The
guard rejects the surrogate when predicted prefix or overhead counters indicate
that the proxy is unreliable.
